\documentclass[11pt,letterpaper]{article}
\usepackage{technical-design-spec}
\setDocTitle{Deployment and Release Readiness --- Template}
\setDocSubtitle{Annotated instructor/authoring instrument}
\setDocVersion{1.0}\setDocStatus{Working instrument / not approved evidence}
\setDocOwner{CI/CD Engineering}\setDocDate{30 October 2026}
\setDocSummary{Use this sheet to record accountable CI/CD engineering decisions; absence of entries is not proof of conformance.}
\begin{document}\makedoctitle\makedocsummary
\section{Control metadata}
\noindent\textbf{Release ID}\\[6pt]\noindent\rule{0.93\linewidth}{0.4pt}\\[14pt]
\noindent\textbf{Source SHA}\\[6pt]\noindent\rule{0.93\linewidth}{0.4pt}\\[14pt]
\noindent\textbf{Artifact digest}\\[6pt]\noindent\rule{0.93\linewidth}{0.4pt}\\[14pt]
\noindent\textbf{Environment}\\[6pt]\noindent\rule{0.93\linewidth}{0.4pt}\\[14pt]
\noindent\textbf{Approver}\\[6pt]\noindent\rule{0.93\linewidth}{0.4pt}\\[14pt]
\section{Release candidate provenance}
Authoring prompt: [Describe the observed state and identify its authoritative evidence.]

Example (illustrative, not operational evidence): The demonstration PR test is required and produces a JUnit XML report; reviewer verifies a deliberately failing test blocks merging.

\noindent\textbf{Evidence / rationale:}\\[8pt]\noindent\rule{0.93\linewidth}{0.4pt}\\[14pt]
\section{Quality and security evidence}
Authoring prompt: [Record the decision boundary and its permitted inputs and outputs.]

Example (illustrative, not operational evidence): The demonstration PR test is required and produces a JUnit XML report; reviewer verifies a deliberately failing test blocks merging.

\noindent\textbf{Evidence / rationale:}\\[8pt]\noindent\rule{0.93\linewidth}{0.4pt}\\[14pt]
\section{Environment approval}
Authoring prompt: [Name objective checks, failure cases and supporting run references.]

Example (illustrative, not operational evidence): The demonstration PR test is required and produces a JUnit XML report; reviewer verifies a deliberately failing test blocks merging.

\noindent\textbf{Evidence / rationale:}\\[8pt]\noindent\rule{0.93\linewidth}{0.4pt}\\[14pt]
\section{Database compatibility}
Authoring prompt: [State who can act, how the action is authorized, and how denial is tested.]

Example (illustrative, not operational evidence): The demonstration PR test is required and produces a JUnit XML report; reviewer verifies a deliberately failing test blocks merging.

\noindent\textbf{Evidence / rationale:}\\[8pt]\noindent\rule{0.93\linewidth}{0.4pt}\\[14pt]
\section{Health checks and abort thresholds}
Authoring prompt: [Give traceable revision IDs, artifacts, risk items and retention policy.]

Example (illustrative, not operational evidence): The demonstration PR test is required and produces a JUnit XML report; reviewer verifies a deliberately failing test blocks merging.

\noindent\textbf{Evidence / rationale:}\\[8pt]\noindent\rule{0.93\linewidth}{0.4pt}\\[14pt]
\section{Rollback and final disposition}
Authoring prompt: [Identify evidence, residual risk, accountable owner and remediation deadline.]

Example (illustrative, not operational evidence): The demonstration PR test is required and produces a JUnit XML report; reviewer verifies a deliberately failing test blocks merging.

\noindent\textbf{Evidence / rationale:}\\[8pt]\noindent\rule{0.93\linewidth}{0.4pt}\\[14pt]
\section{Decision and sign-off}
\noindent\textbf{Final outcome}\\[6pt]\noindent\rule{0.93\linewidth}{0.4pt}\\[14pt]
\noindent\textbf{Residual risks / exception expiry}\\[6pt]\noindent\rule{0.93\linewidth}{0.4pt}\\[14pt]
\noindent\textbf{Owner signature}\\[6pt]\noindent\rule{0.93\linewidth}{0.4pt}\\[14pt]
\noindent\textbf{Independent reviewer}\\[6pt]\noindent\rule{0.93\linewidth}{0.4pt}\\[14pt]
\noindent\textbf{Review date}\\[6pt]\noindent\rule{0.93\linewidth}{0.4pt}\\[14pt]
\end{document}
